% Audited unnumbered theorem and its restricted tetrahedral application.
\paragraph{Endpoint lifting through regular surface families
(audited auxiliary theorem).}
Let $k$ be an algebraically closed field of characteristic $p>0$.
Let $R$ be an excellent henselian local $k$-algebra with residue field
$k$, obtained by strict localization of a finite-type $k$-scheme.
Let $Z\to\operatorname{Spec}R$ be a regular projective threefold,
an isomorphism outside the closed point.  Suppose its reduced closed
fibre $E$ is a simple normal crossing divisor, and that there is an
effective Cartier divisor $D$ supported on all of $E$ such that
$\mathcal O_D(-D)$ is ample.  The scheme $D$ may be nonreduced.
Let $P/k$ be a smooth connected projective curve, and fix two points
$o,t\in P(k)$.  Put $T=D\times P$.

For every $\beta\in F_\gamma^2K_0(T)_{\mathbf Q}$ there is an actual
class $b\in K_0(Z\times P)_{\mathbf Q}$ satisfying
\begin{equation}
 (b|_{Z\times\{u\}})|_E
       =(\beta|_{D\times\{u\}})|_E\quad(u=o,t).
                                                        \tag{EL.1}
\end{equation}
Here $K_0$ is the group of perfect complexes, and $F_\gamma$ is its
vector-bundle gamma filtration.  No equality over the entire family
$E\times P$ is asserted.  This is an endpoint theorem on a fixed
proper parameter curve, not a general specialization theorem for
vector bundles or $K$-theory.

\emph{1. Perfect classes modulo the third gamma filtration.}
The schemes $D$, $E$, and $T$ are projective over $k$, so they have
the resolution property.  To justify this also for the nonreduced
$D$, some power $\mathfrak m^M$ annihilates $\mathcal O_D$; the
Artinian local ring $R/\mathfrak m^M$ is finite-dimensional over
the given coefficient field $k$.  Projectivity over this finite
$k$-algebra gives projectivity over $k$.
Their perfect $K_0$ is the vector-bundle
$K_0$, with its special lambda-ring structure.  For a noetherian
quasi-projective scheme of dimension $d$, the dimension theorem is
\begin{equation}
                         F_\gamma^{d+1}K_0=0.
                                                        \tag{EL.2}
\end{equation}
We use its actual singular-scheme scope, not a regular-scheme
substitute: see \citet[Chapter II, Remark 8.8.5, p.~83]
{D_Weibel2013KbookII}, which cites SGA~6, VI.6.6.
In particular $F_\gamma^3K_0(E)=0$.  A representation of $\beta$
modulo $F_\gamma^3K_0(T)_{\mathbf Q}$ therefore gives its exact
restrictions at both endpoints after further pullback to $E$.

We give a perfect-complex representation retaining classes in the
kernel of the Cartan map.  Modulo $F_\gamma^3$, $F_\gamma^2$ is
generated by $\gamma^2(v)$ and products $v w$, for rank-zero classes
$v,w$.  Every such $v$ can be written $[A]-[B]$, with $A,B$ of the
same rank and sufficiently generated first jets on every smooth
stratum used below.  Indeed, if $L$ is very ample and $V$ is any
vector bundle, choose $m$ with $V^\vee\otimes L^m$ globally generated.
Dualizing a generating surjection and tensoring by $L^m$ gives
\[
 0\longrightarrow V\longrightarrow (L^m)^N
       \longrightarrow Q\otimes L^m\longrightarrow0,
\]
where $Q$ is globally generated.  Taking $m$ still larger supplies
first jets: products of generating sections of $Q$ and first-jet
generators of $L^m$ generate all first jets of $Q\otimes L^m$.
Apply this to finitely many bundles and add the same positive bundle
to both sides if necessary.  Their determinants can likewise be
made very ample.  The universal identities give, modulo $F_\gamma^3$,
\[
 \gamma^2([A]-[B])
 =\gamma^2([A]-n)-\gamma^2([B]-n)
       -(\det A-1)(\det B-1)+(\det B-1)^2.
\]
Here a line bundle denotes its $K_0$ class.  The equality follows
from the special lambda-ring identities and
$([A]-n)\equiv\det A-1\pmod {F_\gamma^2}$; it is valid in every
characteristic.  Products of arbitrary line differences can be
written, modulo $F_\gamma^3$, using differences of positive line
bundles.  Thus it suffices to treat second gamma classes of these
generated bundles and products of two positive line differences.

For a rank-$n$ bundle $A$ choose a general homomorphism
$\mathcal O_T^{n-1}\to A$.  The locus of rank at most $n-3$ has
expected codimension six and is empty.  The rank-$n-2$ locus $C$
has codimension two and is a regular immersion.  Locally an
invertible $(n-2)$-minor reduces the map to an identity block and
a column $(a,b)$; the pair $a,b$ is a regular sequence.  Its Koszul
sequence is the local verification of the global exact sequence
\[
 0\longrightarrow\mathcal O_T^{n-1}\longrightarrow A
      \longrightarrow I_C\otimes\det A\longrightarrow0.
\]
Consequently
\begin{equation}
 [\mathcal O_C]
   =1-[A\otimes(\det A)^{-1}]+(n-1)[(\det A)^{-1}]
   \equiv\gamma^2([A]-n)\pmod {F_\gamma^3}.
                                                        \tag{EL.3}
\end{equation}
The equality is in perfect $K_0(T)$.  It is not merely an equality
of ordinary cycles on the singular threefold $T$.  The congruence
also follows from the exact universal identity
$\det A-[A]+n-1=\sum_{j=2}^n\gamma^j([A]-n)$: divide by
$\det A$, whose difference from $1$ belongs to $F_\gamma^1$.
The identity itself follows after the injective flag-bundle
splitting principle by expanding the product of line roots.
Products of positive line differences are represented modulo
$F_\gamma^3$ by the perfect Koszul classes of general complete
intersections of two sections.  Hence $\beta$, modulo $F_\gamma^3$,
is a finite rational linear combination of such $[\mathcal O_C]$.

The choices can satisfy all the following additional conditions
simultaneously.  The reduced curve meets each component stratum
$E_i\times P$ transversely; it meets each double stratum in finitely
many transverse points; it avoids all triple strata; and its fibres
over $o,t$ are finite regular immersions of codimension two in $D$.
The full curve $C$ need not be flat over every point of $P$;
only these two selected fibres will be used in the endpoint identity.
These are characteristic-free first-jet incidence conditions.
On a smooth stratum of dimension $d$, rank drop by one has
codimension two, and failure of transversality at that rank has
total codimension at least $d+1$.  Deeper rank drop has codimension
six.  First-jet generation makes these the actual incidence
codimensions.  Apply this also to the component and intersection
strata in the two endpoint fibres.  The corresponding bad sets
have images whose closures are proper in the finite-dimensional
parameter space, so their finite union is avoidable over the infinite field
$k$.  The same argument for two sections handles complete
intersections.  Since $D$ is Cartier in a regular scheme, it is
Cohen--Macaulay; transversality on its reduced strata makes the
two local equations a regular sequence also on $D$, with its
given multiplicities.

\emph{2. Put each perfect curve on an actual regular surface family.}
Fix one of these curves $C\subset T$.  Choose a sufficiently high
power $\mathcal A$ of a line bundle on $Z\times P$ ample relative
to $\operatorname{Spec}R$.  It can be taken as an exterior product
of a relatively ample bundle on $Z$ and an ample bundle on $P$.
Serre vanishing makes $I_C\mathcal A|_T$ generated, and makes its
global sections surject onto the sections of
$(I_C/I_C^2)\mathcal A|_C$.  The latter is a rank-two bundle on
the one-dimensional scheme $C$.  A general section is nowhere
zero on $C_{\mathrm{red}}$: the zero condition has codimension two
over a space of dimension one.  Choose such a section, general
also away from $C$, and let $H$ be its zero scheme in $T$.
Then $C$ is an effective Cartier divisor on $H$.  Locally, the
chosen equation is a linear combination of the regular-sequence
generators $a,b$ with one coefficient a unit, so the complementary
generator defines $C$ in $H$.

Choose the section transverse to every reduced normal-crossing
stratum away from $C$.  Along $C$, the first-jet transversality in
Step~1 and the nowhere-zero conormal combination give a nonzero
ambient differential.  This includes the finitely many double
stratum intersections; triple strata were avoided by $C$.
In addition, the section can be required not to vanish identically
on any surface $E_i\times\{v\}$, for every $v\in P$.  For a high
twist the space of its restrictions on each such surface has
dimension at least two, uniformly in $v$: take a generator of
a first twisted ideal that is nonzero on the surface and multiply
it by two independent restrictions from a further very ample
factor.  The condition of being
zero on that surface has codimension at least two; the parameter
$v$ contributes only one dimension.  The incidence union is
therefore a proper subset.  This makes $H\to P$ a projective flat
family of curves, with the indicated multiplicities allowed.
Flatness follows from the fibrewise nonzerodivisor criterion for
the Cartier equation on the Cohen--Macaulay scheme $D\times P$.

Relative Serre vanishing on $Z\times P$ gives the surjection
\[
 H^0(Z\times P,\mathcal A)\longrightarrow H^0(T,\mathcal A|_T)
\]
from $H^1(\mathcal A(-D\times P))=0$.  Lift the chosen equation,
and let $W\subset Z\times P$ be its zero divisor.  The source
$W$ is regular along its closed fibre: this is exactly the
nonzero ambient differential just imposed on every reduced
stratum of $E\times P$.  It is regular everywhere.  Indeed its
nonregular locus is closed, since the schemes are excellent, and
its proper image in the local scheme $\operatorname{Spec}R$ would
contain the closed point if it were nonempty.  Its intersection
with that fibre is empty.  The same fibrewise nonzerodivisor
criterion gives $W\to P$ flat: on every parameter fibre its
equation is nonzero on every regular component of $Z$.  Every
such component meets the closed fibre by properness, and the
equation was required not to vanish identically on any of its
exceptional surface components.  Thus it is a nonzerodivisor on
every fibre of the flat morphism $Z\times P\to P$.
We have constructed the actual
fixed subscheme
\[
 W\subset Z\times P,\qquad H=W\cap(D\times P),
\]
proper over $R$, with $C$ Cartier on $H$.  All finite first-jet
and section choices may equivalently be made after descending to
one pointed etale finite-type model of $Z$ and then pulling back;
no Bertini statement for an unspecified non-finite-type variety
is being assumed.  The only use of relative Serre vanishing here
is its usual projective form
\citep[Lemma 30.17.1]{frC3azStacksSerre}.

\emph{3. One finite power lifts the whole line bundle on $H$.}
Put $J=\mathcal O_W(-D\times P)$ and
$L=J|_H$.  This is ample on every fibre of the projective morphism
$H\to P$, since $\mathcal O_D(-D)$ is ample.  Relative Serre
vanishing therefore gives $R^i(H\to P)_*L^n=0$ for $i>0$ and all
sufficiently large $n$.  As $P$ is a curve, the Leray sequence gives
\begin{equation}
                       H^2(H,L^n)=0\quad(n\gg0).
                                                        \tag{EL.4}
\end{equation}
We do not assert that $H^1(H,L^n)$ vanishes.

Start with $Q_0=\mathcal O_H(-C)$.  For the thickenings
$H_n=V(J^n)\subset W$, the obstruction to extending a line bundle
from $H_n$ to $H_{n+1}$ lies in $H^2(H,L^n)$, by the square-zero
unit sequence\citep[Lemma 37.4.1]{frC3azStacksPicThickening}.
These groups are killed by $p$.  At each of the finitely many
steps before (EL.4) applies, raise the current partial bundle and
all its preceding restrictions to the $p$th tensor power before
extending.  After a finite number $e$ of such steps, there is a
compatible system on every $H_n$ restricting to $Q_0^{p^e}$ on
$H$.  Subsequent obstruction groups vanish, so the exponent never
increases again.  No uniform Frobenius bound for arbitrary vector
bundles is used.

The ideals $J$ and $\mathfrak m\mathcal O_W$ have the same radical
and cofinal powers.  Projective Grothendieck existence algebraizes
this formal line bundle on $W_{\widehat R}$
\citep[Lemma 30.24.3]{frC3azStacksExistence}; algebraizing its inverse
and using full faithfulness verifies invertibility.  The regular
completion map $R\to\widehat R$ is a filtered colimit of smooth
$R$-algebras by Popescu
\citep[Lemma 15.51.14; Theorem 16.12.1]
{frC3azStacksCompletion,frC3azStacksPopescu}.
The algebraized bundle descends to one such algebra by finite
presentation\citep[Lemma 32.10.3(2)]{frC3azStacksLimitBundles}.
Choose $M$ with $\mathfrak m^M\mathcal O_W\subset J$.
The descended algebra's map to $\widehat R/\mathfrak m^M=R/\mathfrak m^M$
lifts to $R$, since $(R,\mathfrak m^M)$ is a henselian pair and the
algebra is smooth\citep[Lemma 15.13.3]{frC3azStacksHenselSmooth}.
The resulting actual line bundle $Q$ on the fixed $W$ satisfies
\begin{equation}
                            Q|_H\simeq Q_0^{p^e}.
                                                        \tag{EL.5}
\end{equation}
Matching this finite quotient preserves the whole proper parameter
curve and the entire first thickening $H$, not just two pointwise
Picard classes.

\emph{4. Exact endpoint equality, without a filtration pushforward.}
Let $j:W\hookrightarrow Z\times P$.  This is a regular Cartier
immersion, so
\[
                         b_C=j_*(1-[Q])
\]
is an actual perfect class.  Fix $u=o$ or $t$.  All intersections
with the parameter fibre and with $D$ are Tor independent by the
flatness and regular-sequence conditions already imposed.  Thus
restriction to $D\times\{u\}$ is
\[
       (j_{H_u})_*\bigl(1-[\mathcal O_{H_u}(-p^e C_u)]\bigr).
\]
The scheme $C_u$ is finite over $k$.  Every line bundle on a finite
Artinian $k$-scheme is trivial.  The Cartier sequence and the
geometric-series identity therefore give, in perfect $K_0(D)$,
\begin{align*}
 (j_{H_u})_*(1-[Q_{0,u}]^{p^e})
 &=\sum_{a=0}^{p^e-1}(j_{H_u})_*
                   \bigl([\mathcal O_{C_u}]\,[Q_{0,u}]^a\bigr)\\
 &=p^e[\mathcal O_{C_u}].
\end{align*}
This is exact; no compatibility of a singular pushforward with
the gamma filtration has been assumed.  Pull back further to $E$.
Dividing $b_C$ by $p^e$ and adding the finite rational combination
of Step~1 gives (EL.1), since the omitted $F_\gamma^3$ part vanishes
on both copies of $E$ by (EL.2).  This proves the endpoint
theorem.

\paragraph{Application to the ordered tetrahedral germ
(proved application).}
For the exact TG construction, TS gives a fixed natural summand
$\kappa K_2(P)_{\mathbf Q}$ on $E\times P$, where
$\kappa\in K_{-2}(E_0)_{\mathbf Q}$ has weight zero.  Let
$\xi\in K_2(P)$ and put $a=\kappa\xi\in K_0(E\times P)_{\mathbf Q}$.
The class $a$ has rank zero, since $\kappa$ restricts to zero in
negative $K$-theory at every point.  Lift $a$ to
$K_0(D\times P)_{\mathbf Q}$ by nilpotent invariance
after inverting $p$\citep[Theorem A]{rootGeisserHesselholt2011}.
For either endpoint $u$, its restriction $a_u$ is square-zero and
has Adams weight two: $\kappa^2=0$ by TS over $k_0$, and field
$K_2$ has weight two.  Choose a positive integer $N$ such that the
lift $Na$ is integral, necessarily of rank zero.  In the special lambda-ring,
\[
 \gamma^2(Na_u)
   =\frac{(Na_u)^2+2Na_u-\psi^2(Na_u)}2=-Na_u.
\]
Apply (EL.1) to $-\gamma^2(Na)/N$.  It yields an actual perfect
family on $Z\times P$ whose endpoint restrictions to $E$ are
exactly $a_o,a_t$.  No global weight decomposition for $K_2(P)$
has been used.

By FF.4, restriction $K_0(Z)_{\mathbf Q}\to K_0(E)_{\mathbf Q}$
is an isomorphism.  Its endpoint lifts are therefore precisely
the unique weight-two lifts used in TC.  The second Chern-character
component of the actual family, defined on finite-type smooth
etale models and then pulled back, proves that their difference
is ordinarily algebraically trivial in $\mathrm{CH}^2(Z)_{\mathbf Q}$.
Ordinary here means an actual fixed product with the smooth proper
curve $P$, as required.

Every field symbol $\{a,b\}\in K_2(k)$ is such an endpoint
difference on a smooth proper curve, by the explicitly proved
proper-curve symbol lemma retained in the top-quotient research:
normalize a component of $(f-1)^2(g-1)^2=c f g$, use the unramified
symbol $\{f,g\}$, and compare $(a,b)$ with $(0,1)$.
The exact localization and Gersten inputs are
\citet[Section 7, Proposition 3.2 and equation (3.3), p.~127;
Theorem 5.11, pp.~133--134]{rootQuillen1973}.
Finite sums of these symbols generate $K_2(k)$.  TC identifies
$\mathrm{CH}^2(X)_{\mathbf Q}$ with precisely this coefficient
group and kills the four degree coordinates by exceptional curves.
Consequently
\[
              B^2(X)_{\mathbf Q}=0
\]
for this specified tetrahedral family in every positive
characteristic and over every algebraically closed extension of
$\overline{\mathbf F}_p$.  The endpoint theorem and this application have full independent
root and two-peer proof audits.  The conclusion is specific to
this tetrahedral family; it does not assert the same result for
every isolated threefold singularity.
